\documentclass[aspectratio=169]{beamer}

\usetheme{Madrid}
\usecolortheme{default}

\title{Project Plan: Synthetic Clinical Note Generation for GLP-1 RA Discontinuation Reason Classification}
\author{Adiba, Shahidul \\ BioNLP Research Group}
\date{\today}

\begin{document}

\begin{frame}
    \titlepage
\end{frame}

\begin{frame}{Plan Overview}
    \tableofcontents
\end{frame}

\section{Motivation and Scope}

\begin{frame}{Problem and Motivation}
    \begin{itemize}
        \item GLP-1 receptor agonists are widely used for diabetes and weight management, but discontinuation is common and clinically important.
        \item Reasons for discontinuation are often buried in free-text clinical notes rather than structured fields.
        \item Supervised NLP models need high-quality labeled data, but expert annotation of real clinical notes is expensive and difficult to scale.
        \item This project focuses on synthetic note generation as a practical path toward building training data for discontinuation reason classification.
    \end{itemize}
\end{frame}

\begin{frame}{Project Objective}
    \begin{itemize}
        \item Build a publishable end-to-end pipeline for synthetic GLP-1 discontinuation note generation, validation, and benchmarking.
        \item Create a diverse synthetic dataset in VA-style EHR format.
        \item Validate note quality through embedding-based diversity checks, LLM review, and human review.
        \item Benchmark the curated dataset on downstream classification tasks.
        \item Document limitations, ethics considerations, and reproducibility practices expected for a dataset paper.
    \end{itemize}
\end{frame}

\begin{frame}{Taxonomy and Planned Scope}
    \textbf{Initial label set}
    \begin{itemize}
        \item Adverse Effects / Tolerability
        \item Access / Cost / System Barriers
        \item Patient Preference / Beliefs
        \item Planned / Goal-Based Discontinuation
    \end{itemize}

    \vspace{0.5em}
    \textbf{Planned expansion areas}
    \begin{itemize}
        \item Main Label
        \item Clinical Decision / Medical Appropriateness
        \item Efficacy / Expectation Mismatch
        \item Other / Unknown
    \end{itemize}
\end{frame}

\begin{frame}{Research Questions and Planned Contributions}
    \small
    \textbf{Research questions}
    \begin{itemize}
        \item Can synthetic GLP-1 discontinuation notes achieve strong realism, diversity, and label fidelity?
        \item How well do LLM-based judgments align with expert human review?
        \item Does a curated synthetic dataset improve downstream classification across prompting and supervised settings?
        \item Which curation steps matter most for downstream performance?
    \end{itemize}

    \textbf{Expected contributions}
    \begin{itemize}
        \item A clinically motivated synthetic dataset and benchmark
        \item A validation pipeline that combines embeddings, LLM judges, and human judges
        \item A multi-model benchmark for discontinuation reason classification
        \item An empirical analysis of when synthetic data helps and where it fails
    \end{itemize}
\end{frame}

\section{Methodology}

\begin{frame}{Planned Synthetic Data Pipeline}
    \begin{enumerate}
        \item \textbf{Generation:} Use few-shot prompting to produce VA-style clinical notes for each discontinuation category.
        \item \textbf{Diversity control:} Apply embedding-based similarity checks to reduce near-duplicate notes.
        \item \textbf{Automated validation:} Use an LLM judge to score realism, alignment, coherence, variability, and GLP-1 relevance.
        \item \textbf{Human validation:} Review a sampled subset with domain-informed annotation criteria.
        \item \textbf{Curation:} Filter low-quality notes and build stable development and test partitions.
    \end{enumerate}
\end{frame}

\begin{frame}{Planned Evaluation Strategy}
    \begin{itemize}
        \item Generate and curate a synthetic dataset covering the target discontinuation categories.
        \item Report descriptive dataset characteristics such as class balance, note length patterns, lexical diversity, and embedding diversity.
        \item Maintain a fixed held-out test set that remains untouched during development.
        \item Measure agreement between automated judgments and human review.
        \item Run ablations to test the impact of diversity filtering, quality filtering, and prompt design choices.
    \end{itemize}
\end{frame}

\begin{frame}{Planned Benchmarking Strategy}
    \small
    \textbf{Model families}
    \begin{itemize}
        \item Qwen family
        \item Gemma family
        \item Llama family
        \item Clinical encoder baselines such as BioClinicalBERT or ClinicalBERT
    \end{itemize}

    \textbf{Evaluation settings}
    \begin{itemize}
        \item Zero-shot prompting
        \item Few-shot prompting
        \item Supervised fine-tuning
        \item Preference-based optimization where appropriate
    \end{itemize}
\end{frame}

\begin{frame}{Validation and Publication Requirements}
    \small
    \begin{itemize}
        \item Define a human evaluation protocol with reviewer expertise, annotation rubric, and agreement analysis.
        \item Include strong discriminative baselines in addition to generative LLMs.
        \item Study the effect of diversity filtering, quality filtering, prompt design, and category granularity through ablations.
        \item Compare automated review decisions with human review outcomes.
        \item Seek external validation against real or expert-authored notes when feasible.
        \item Prepare dataset documentation, ethics notes, bias analysis, and reproducibility details for release.
    \end{itemize}
\end{frame}

\section{Roadmap}

\begin{frame}{Next Steps}
    \begin{itemize}
        \item Run the full validation pipeline on the expanded synthetic corpus.
        \item Apply diversity filtering, quality filtering, and structured error analysis.
        \item Finalize stable train, development, and test splits.
        \item Benchmark downstream classification across prompting and supervised settings.
        \item Compare LLM-judge decisions with human review outcomes.
        \item Expand beyond the initial label set using the broader discontinuation taxonomy.
        \item Compare the resulting dataset against real-world EHR-style material for external validity.
    \end{itemize}
\end{frame}

\begin{frame}{Closing Plan Summary}
    \begin{block}{Project direction}
        The study will motivate the task, define a clinically grounded taxonomy, build a synthetic note generation and curation pipeline, validate the resulting data with automated and human review, and benchmark its value for downstream classification.
    \end{block}

    \vspace{0.5em}
    \begin{itemize}
        \item Focus on a clear and defensible dataset paper narrative
        \item Emphasize validation, benchmarking, and reproducibility
        \item Keep the presentation centered on planned work rather than interim results
    \end{itemize}
\end{frame}

\end{document}
